\scanpage{71}
\mspage{\deleted{\unsure{69|79}}\unsure{70|79}}

Il revient au $=$ de dire que $F$ et $F'$ ont $=$ cosupport, ou encore, même support.

\unsure{Donnons nous} un support.

% Two short columns of marks in the left margin opposite this line.
\margin{$\unsure{!|\vdots}\qquad\unsure{!|\vdots}$}
Une autre façon de dire que $F,G\ \text{disj.}\Longrightarrow F',G\ \text{disj}$, est ainsi :

\uline{Corollaire} Le cosupport \inserted{dans $\mathcal F$} de tout objet $G$ de $\mathcal F$ est fermé, non \unsure{seult|seulemt} pour $\leq$, mais aussi pour $\unsure{\preceq|\leqslant}$ -- c'est donc ce qu'on appelle une \uline{sous}-\unsure{contrée|contrie}.

Donc \unsure{itou} pour toute partie $\Phi\subset\mathcal F$.

En d'autres termes, \deleted{toute} partie de $\mathcal F$ qui est un $\mathcal F$-support (ou ``$\mathcal F$-saturée, sous-entendu : en termes de $\leq$\unsure{,)|,) ;} cf page $67$) est \unsure{automatiquemt|automatiquement} \deleted{\unsure{f}}une sous-contrée.

% Cancelled draft, written upside down. The diagonal strokes run across the whole block.
\mspage{\deleted{\unsure{5\deleted{\unsure{4|3}}5}}}

\deleted{un ensemble, et $\deleted{\unsure{\mathcal F|F}}\Phi$ \unsure{un} ens \unsure{de|des} parties de $M$, satisfaisant la condition précédente \unsure{(au --)|(ou --)|(en --)}.}

\deleted{Alors $\Phi$, muni de la relation d'inclusion, satisfait $C\,1$, et ses objets irréductibles sont les $A_x$ $(x\in M)$, l'application $x\unsure{\mathrel{\vdash\joinrel\dashv}|\mapsto} A_x$ étant une bijection de $M$ sur l'ens des \deleted{\unsure{ses|son} \unsure{par|p}} objets irréductibles de $\Phi$.}

\deleted{\unsure{Dém|Dire|Bien} $A_x$ est irréductible, car si $A_x=\bigcup_{i\in I} A_i$}
